% !Mode:: "TeX:UTF-8"
% !TEX root = ../../main.tex
% ============================================================
% TOPIC 5
% ============================================================
\mashgulot{5}{The inverse matrix}
\vspace{3pt}
\begin{tarif}
Let $A$ be a square matrix of order $n$. A matrix $A^{-1}$ such that
\[\kf{A\cdot A^{-1} = A^{-1}\cdot A = E,}\]
where $E$ is the identity matrix of order $n$, is called the
\textbf{inverse} of $A$.
\end{tarif}
Not every square matrix has an inverse. The answer to the question of
existence is given by the determinant.
\vspace{3pt}
\begin{teorema}
A matrix $A$ has an inverse if and only if $\det A \neq 0$, that is, if and
only if $A$ is \textbf{nonsingular}. In this case the inverse is unique.
\end{teorema}
\emph{The inverse via cofactors}
\vspace{3pt}
\begin{tarif}
If $A$ is an $n\times n$ matrix with $\det A \neq 0$, then
\[\kf{A^{-1} = \frac{1}{\det A}\cdot \widetilde{A}^{T},}\]
where $\widetilde{A} = (A_{ij})$ is the matrix of cofactors (algebraic
complements) of $A$. Its transpose $\widetilde{A}^{T}$ is called the
\textbf{adjugate} (or adjoint) matrix and is denoted by $\operatorname{adj}A$.
\end{tarif}
For $2\times 2$ and $3\times 3$ matrices the computation is organized as
follows.
\begin{enumerate}[label=\arabic*.]
\item Compute $\det A$. If $\det A = 0$, the inverse does not exist.
\item For every entry $a_{ij}$ find the cofactor $A_{ij} = (-1)^{i+j}M_{ij}$,
where $M_{ij}$ is the corresponding minor.
\item Form the cofactor matrix $\widetilde{A}=(A_{ij})$ and transpose it to
obtain $\operatorname{adj}A=\widetilde{A}^{T}$.
\item Divide by the determinant: $A^{-1} = \dfrac{1}{\det A}\cdot \widetilde{A}^{T}$.
\end{enumerate}
It is always a good idea to finish by checking that $A\cdot A^{-1} = E$.

\textbf{The inverse of a $2\times 2$ matrix.}
If $A=\begin{pmatrix}a&b\\c&d\end{pmatrix}$ and $\det A = ad-bc\neq 0$, then
\[\kf{A^{-1} = \frac{1}{ad-bc}\begin{pmatrix}d & -b\\ -c & a\end{pmatrix}.}\]
In words: interchange the diagonal entries, change the signs of the
off-diagonal ones and divide by the determinant.
\begin{misol}
Find the inverse of $A=\begin{pmatrix}3 & 1\\ 5 & 2\end{pmatrix}$.
\begin{yechish}
Here $\det A = 3\cdot 2 - 1\cdot 5 = 6-5 = 1 \neq 0$, so the inverse exists and
\[\kf{A^{-1} = \frac{1}{1}\begin{pmatrix}2 & -1\\ -5 & 3\end{pmatrix}
= \begin{pmatrix}2 & -1\\ -5 & 3\end{pmatrix}.}\]
As a check,
\[\kf{A\cdot A^{-1} = \begin{pmatrix}3&1\\5&2\end{pmatrix}\begin{pmatrix}2&-1\\-5&3\end{pmatrix}}\]
\[\kf{= \begin{pmatrix}6-5 & -3+3\\ 10-10 & -5+6\end{pmatrix}
= \begin{pmatrix}1&0\\0&1\end{pmatrix} = E.}\]
\end{yechish}
\end{misol}
\begin{misol}
For $A=\begin{pmatrix}4 & -2\\ 6 & -3\end{pmatrix}$ we find
$\det A = 4\cdot(-3) - (-2)\cdot 6 = -12+12 = 0$. The determinant vanishes,
so $A$ has \textbf{no inverse}.
\end{misol}
\textbf{The inverse of a $3\times 3$ matrix.}
\vspace{3pt}
\begin{misol}
Find the inverse of
$A=\begin{pmatrix}2 & 5 & 7\\ 6 & 3 & 4\\ 5 & -2 & -3\end{pmatrix}$.
\begin{yechish}
We first expand the determinant along the first row:
\begin{align*}
\det A &= 2\bigl(3 \cdot (-3) - 4 \cdot (-2)\bigr) - 5\bigl(6 \cdot (-3) - 4 \cdot 5\bigr)\\
&\quad + 7\bigl(6 \cdot (-2) - 3 \cdot 5\bigr)\\
&= 2\cdot(-1) - 5\cdot(-38) + 7\cdot(-27)\\
&= -2 + 190 - 189 = -1 \neq 0.
\end{align*}
Next we compute the nine cofactors:
\begin{align*}
A_{11} &= +\begin{vmatrix}3&4\\-2&-3\end{vmatrix} = -9+8 = -1, \\
A_{12} &= -\begin{vmatrix}6&4\\5&-3\end{vmatrix} = -(-18-20) = 38, \\
A_{13} &= +\begin{vmatrix}6&3\\5&-2\end{vmatrix} = -12-15 = -27, \\
A_{21} &= -\begin{vmatrix}5&7\\-2&-3\end{vmatrix} = -(-15+14) = 1, \\
A_{22} &= +\begin{vmatrix}2&7\\5&-3\end{vmatrix} = -6-35 = -41,
\end{align*}
\begin{align*}
A_{23} &= -\begin{vmatrix}2&5\\5&-2\end{vmatrix} = -(-4-25) = 29, \\
A_{31} &= +\begin{vmatrix}5&7\\3&4\end{vmatrix} = 20-21 = -1, \\
A_{32} &= -\begin{vmatrix}2&7\\6&4\end{vmatrix} = -(8-42) = 34, \\
A_{33} &= +\begin{vmatrix}2&5\\6&3\end{vmatrix} = 6-30 = -24.
\end{align*}
Writing the cofactors of the $i$th row of $A$ into the $i$th column, we
obtain the adjugate matrix
\[\kf{\operatorname{adj}A=\widetilde{A}^{T} = \begin{pmatrix}-1&1&-1\\38&-41&34\\-27&29&-24\end{pmatrix},}\]
and finally
\[\kf{A^{-1} = \frac{1}{-1}\begin{pmatrix}-1&1&-1\\38&-41&34\\-27&29&-24\end{pmatrix}
= \begin{pmatrix}1&-1&1\\-38&41&-34\\27&-29&24\end{pmatrix}.}\]
The reader is encouraged to verify that $A\cdot A^{-1}=E$; the computation is
also repeated by computer at the end of this section.
\end{yechish}
\end{misol}
\emph{Properties of the inverse}

Let $A$ and $B$ be nonsingular $n\times n$ matrices and let $k\neq 0$ be a
scalar. Then
\begin{enumerate}[label=\arabic*)]
\item $(A^{-1})^{-1} = A$;
\item $(AB)^{-1} = B^{-1}A^{-1}$ \hfill (the order of the factors is reversed);
\item $(A^{T})^{-1} = (A^{-1})^{T}$ \hfill (inversion commutes with transposition);
\item $(kA)^{-1} = \dfrac{1}{k}A^{-1}$;
\item $\det(A^{-1}) = \dfrac{1}{\det A}$;
\item the inverse $A^{-1}$ is unique.
\end{enumerate}
\vspace{3pt}
\begin{misol}
Let us illustrate property 2 with $A=\begin{pmatrix}1&2\\0&1\end{pmatrix}$ and
$B=\begin{pmatrix}1&0\\3&1\end{pmatrix}$. By the $2\times 2$ formula,
$A^{-1}=\begin{pmatrix}1&-2\\0&1\end{pmatrix}$ and
$B^{-1}=\begin{pmatrix}1&0\\-3&1\end{pmatrix}$. Further,
\[\kf{AB = \begin{pmatrix}1&2\\0&1\end{pmatrix}\begin{pmatrix}1&0\\3&1\end{pmatrix}
= \begin{pmatrix}7&2\\3&1\end{pmatrix},}\]
\[\kf{(AB)^{-1} = \frac{1}{7-6}\begin{pmatrix}1&-2\\-3&7\end{pmatrix}
= \begin{pmatrix}1&-2\\-3&7\end{pmatrix},}\]
and indeed
\[\kf{B^{-1}A^{-1} = \begin{pmatrix}1&0\\-3&1\end{pmatrix}\begin{pmatrix}1&-2\\0&1\end{pmatrix}
= \begin{pmatrix}1&-2\\-3&7\end{pmatrix}.}\]
\end{misol}
\textbf{Matrix equations.}
The inverse matrix allows us to solve matrix equations of the form $AX=B$,
$XA=B$ and $AXB=C$ with nonsingular $A$ (and $B$ in the last case).
\vspace{3pt}
\begin{tarif}
\begin{itemize}[leftmargin=1.3em]
\item $AX=B \;\Longrightarrow\; X = A^{-1}B$ (multiply by $A^{-1}$ on the left).
\item $XA=B \;\Longrightarrow\; X = BA^{-1}$ (multiply by $A^{-1}$ on the right).
\item $AXB=C \;\Longrightarrow\; X = A^{-1}CB^{-1}$.
\end{itemize}
\end{tarif}
Since matrix multiplication is not commutative, the side on which we
multiply matters: in general $A^{-1}B\neq BA^{-1}$.
\begin{misol}
Solve the equation $AX = B$, where
$A=\begin{pmatrix}3&1\\5&2\end{pmatrix}$ and $B=\begin{pmatrix}1&0\\-1&2\end{pmatrix}$.
\begin{yechish}
The inverse of $A$ was found in the first example of this section:
$A^{-1}=\begin{pmatrix}2&-1\\-5&3\end{pmatrix}$. Hence
\[\kf{X = A^{-1}B = \begin{pmatrix}2&-1\\-5&3\end{pmatrix}\begin{pmatrix}1&0\\-1&2\end{pmatrix}}\]
\[\kf{= \begin{pmatrix}2+1 & 0-2\\ -5-3 & 0+6\end{pmatrix}
= \begin{pmatrix}3&-2\\-8&6\end{pmatrix}.}\]
Substituting back, we get
\[\kf{AX = \begin{pmatrix}3&1\\5&2\end{pmatrix}\begin{pmatrix}3&-2\\-8&6\end{pmatrix}
= \begin{pmatrix}9-8&-6+6\\15-16&-10+12\end{pmatrix}
= \begin{pmatrix}1&0\\-1&2\end{pmatrix} = B.}\]
\end{yechish}
\end{misol}
\begin{misol}
Solve the equation $XA = B$, where
$A=\begin{pmatrix}1&2\\3&5\end{pmatrix}$ and $B=\begin{pmatrix}4&-1\\2&3\end{pmatrix}$.
\begin{yechish}
Since $\det A = 5-6=-1\neq 0$, the matrix $A$ is invertible and
$A^{-1}=\frac{1}{-1}\begin{pmatrix}5&-2\\-3&1\end{pmatrix}
=\begin{pmatrix}-5&2\\3&-1\end{pmatrix}$. Now we multiply by $A^{-1}$ on the
right:
\[\kf{X = BA^{-1} = \begin{pmatrix}4&-1\\2&3\end{pmatrix}\begin{pmatrix}-5&2\\3&-1\end{pmatrix}}\]
\[\kf{= \begin{pmatrix}-20-3&8+1\\-10+9&4-3\end{pmatrix}
= \begin{pmatrix}-23&9\\-1&1\end{pmatrix}.}\]
\end{yechish}
\end{misol}
\subsection{The inverse matrix in information security}
\textbf{Decryption in the Hill cipher.}
In the Hill cipher a plaintext block $p$ is encrypted by the rule
$c = Kp \pmod{26}$. To decrypt, we need the inverse key:
\[\kf{p = K^{-1}c \pmod{26}.}\]
Here $K^{-1}$ must be computed modulo $26$. We find $\det K \pmod{26}$ and
check that $\gcd(\det K, 26) = 1$; then we find the multiplicative inverse
$(\det K)^{-1} \pmod{26}$ and put
\[\kf{K^{-1} \equiv (\det K)^{-1}\cdot \operatorname{adj}K \pmod{26}.}\]
\vspace{3pt}
\begin{misol}
Take the key matrix $K=\begin{pmatrix}1&2\\3&5\end{pmatrix}$. Then
$\det K = 5-6=-1\equiv 25 \pmod{26}$. Since $25\cdot 25 = 625 = 24\cdot 26 + 1$,
we have $25^{-1}\equiv 25\pmod{26}$. Moreover,
$\operatorname{adj}K = \begin{pmatrix}5&-2\\-3&1\end{pmatrix}
\equiv \begin{pmatrix}5&24\\23&1\end{pmatrix}\pmod{26}$, and therefore
\[\kf{K^{-1}\equiv 25\cdot\begin{pmatrix}5&24\\23&1\end{pmatrix}
= \begin{pmatrix}125&600\\575&25\end{pmatrix}
\equiv \begin{pmatrix}21&2\\3&25\end{pmatrix}\pmod{26}.}\]
If the ciphertext is $c=\begin{pmatrix}10\\3\end{pmatrix}$, then
$p = K^{-1}c = \begin{pmatrix}21&2\\3&25\end{pmatrix}\begin{pmatrix}10\\3\end{pmatrix}
= \begin{pmatrix}216\\105\end{pmatrix}\equiv\begin{pmatrix}8\\1\end{pmatrix}\pmod{26}$.
With the encoding $A=0$, $B=1$, \dots, $Z=25$ these numbers are the letters
“I” and “B”.
\end{misol}
\textbf{Digital signatures and authentication.}
Some challenge–response protocols use matrix transformations: the server
sends a matrix $A$, and the user proves his or her identity by returning its
inverse $A^{-1}$ (or the result of applying it). How fast the inverse can be
computed largely determines the efficiency of such protocols.

\textbf{Decoding error-correcting codes.}
In coding theory, a received word is decoded with the help of the inverse of
a suitable part of the generator matrix. For instance, to extract the
information bits from a codeword of a systematic code one applies the
inverse of the information part of $G$.

\textbf{Transformations in computer graphics.}
Visualization systems, including those that display encrypted 3D data,
describe rotations, scalings and translations by matrices. To undo a
transformation, one applies the inverse matrix.

\subsection{Computing with Python, Maple and Mathcad}
We repeat by computer the inversion of the $3\times 3$ matrix
$A=\begin{pmatrix}2 & 5 & 7\\ 6 & 3 & 4\\ 5 & -2 & -3\end{pmatrix}$ and the
solution of the matrix equation $XA=B$ from the last example (there the
matrices are called \texttt{P} and \texttt{Q}).
\begin{Verbatim}[fontsize=\footnotesize,frame=single,label={\small Python}]
from sympy import Matrix, eye
A = Matrix([[2, 5, 7], [6, 3, 4], [5, -2, -3]])
print("det A =", A.det())
Ainv = A.inv()
for row in Ainv.tolist():
    print(row)
print(A * Ainv == eye(3))
P = Matrix([[1, 2], [3, 5]])
Q = Matrix([[4, -1], [2, 3]])
print(Q * P.inv())      # solution of X*P = Q
\end{Verbatim}
\begin{Verbatim}[fontsize=\footnotesize,frame=single,label={\small Output}]
det A = -1
[1, -1, 1]
[-38, 41, -34]
[27, -29, 24]
True
Matrix([[-23, 9], [-1, 1]])
\end{Verbatim}
The method \texttt{adjugate()} returns $\operatorname{adj}A$ if one wants to
follow the cofactor method literally. In Maple:
\begin{Verbatim}[fontsize=\footnotesize,frame=single,label={\small Maple}]
with(LinearAlgebra):
A := Matrix([[2,5,7], [6,3,4], [5,-2,-3]]):
Determinant(A);             # -1
Adjoint(A);                 # adjugate matrix
Ainv := MatrixInverse(A);
A . Ainv;                   # identity matrix
P := Matrix([[1,2], [3,5]]):
Q := Matrix([[4,-1], [2,3]]):
X := Q . MatrixInverse(P);  # [[-23,9],[-1,1]]
\end{Verbatim}
In Mathcad Prime the inverse is obtained with the power $-1$ and the
determinant with the operator $|A|$; the identity matrix is
\texttt{identity(3)}.
\begin{Verbatim}[fontsize=\footnotesize,frame=single,label={\small Mathcad}]
A := | 2   5   7 |
     | 6   3   4 |
     | 5  -2  -3 |
|A| = -1
         |   1   -1    1 |
A^-1 =   | -38   41  -34 |
         |  27  -29   24 |
          | 1  0  0 |
A*A^-1 =  | 0  1  0 |
          | 0  0  1 |
P := | 1  2 |     Q := | 4  -1 |
     | 3  5 |          | 2   3 |
X := Q*P^-1       X = | -23  9 |
                      |  -1  1 |
\end{Verbatim}

\subsection{Exercises on the inverse matrix}
\bigskip
\emph{1. Find the inverse of each $2\times 2$ matrix, or show that it does not exist.}
\bigskip
\begin{multicols}{3}\raggedright
\begin{enumerate}[label=\arabic*.,  leftmargin=*]
\item $\begin{pmatrix}3&1\\5&2\end{pmatrix}$
\item $\begin{pmatrix}4&-2\\6&-3\end{pmatrix}$
\item $\begin{pmatrix}1&2\\3&5\end{pmatrix}$
\item $\begin{pmatrix}-2&3\\1&-1\end{pmatrix}$
\item $\begin{pmatrix}5&0\\0&-3\end{pmatrix}$
\item $\begin{pmatrix}2&4\\1&2\end{pmatrix}$
\item $\begin{pmatrix}-1&3\\2&-5\end{pmatrix}$
\item $\begin{pmatrix}7&2\\3&1\end{pmatrix}$
\item $\begin{pmatrix}6&-4\\-9&6\end{pmatrix}$
\item $\begin{pmatrix}1&-1\\-1&2\end{pmatrix}$
\item $\begin{pmatrix}0&1\\1&0\end{pmatrix}$
\item $\begin{pmatrix}3&6\\2&4\end{pmatrix}$
\item $\begin{pmatrix}-4&1\\2&3\end{pmatrix}$
\item $\begin{pmatrix}1&0\\0&1\end{pmatrix}$
\item $\begin{pmatrix}2&-3\\-4&6\end{pmatrix}$
\item $\begin{pmatrix}5&2\\-1&3\end{pmatrix}$
\item $\begin{pmatrix}-3&-1\\6&2\end{pmatrix}$
\item $\begin{pmatrix}8&5\\3&2\end{pmatrix}$
\item $\begin{pmatrix}1&4\\2&7\end{pmatrix}$
\item $\begin{pmatrix}-2&0\\0&-5\end{pmatrix}$
\item $\begin{pmatrix}4&3\\-8&-6\end{pmatrix}$
\item $\begin{pmatrix}9&-1\\4&1\end{pmatrix}$
\item $\begin{pmatrix}-5&2\\10&-4\end{pmatrix}$
\item $\begin{pmatrix}2&1\\5&3\end{pmatrix}$
\item $\begin{pmatrix}3&-5\\1&-2\end{pmatrix}$
\item $\begin{pmatrix}0&-2\\3&0\end{pmatrix}$
\item $\begin{pmatrix}-1&-4\\3&11\end{pmatrix}$
\item $\begin{pmatrix}6&1\\-2&3\end{pmatrix}$
\item $\begin{pmatrix}10&5\\-2&-1\end{pmatrix}$
\item $\begin{pmatrix}7&-3\\2&1\end{pmatrix}$
\end{enumerate}
\end{multicols}
\emph{2. Find the inverse of each $3\times 3$ matrix by the cofactor method.}
\bigskip
\begin{multicols}{3}\raggedright
\begin{enumerate}[label=\arabic*.,  leftmargin=*]
\item $\begin{pmatrix}2&5&7\\6&3&4\\5&-2&-3\end{pmatrix}$
\item $\begin{pmatrix}1&-1&2\\0&3&-1\\2&1&0\end{pmatrix}$
\item $\begin{pmatrix}3&0&1\\-1&2&-4\\0&5&3\end{pmatrix}$
\item $\begin{pmatrix}-2&1&3\\4&0&-1\\1&-2&5\end{pmatrix}$
\item $\begin{pmatrix}1&2&0\\0&1&-3\\-1&0&4\end{pmatrix}$
\item $\begin{pmatrix}4&-1&2\\3&5&0\\-2&0&1\end{pmatrix}$
\item $\begin{pmatrix}0&1&-2\\3&0&4\\-1&2&0\end{pmatrix}$
\item $\begin{pmatrix}1&0&3\\-2&1&0\\4&-3&2\end{pmatrix}$
\item $\begin{pmatrix}-3&2&1\\0&-4&5\\2&0&-1\end{pmatrix}$
\item $\begin{pmatrix}5&-2&0\\1&3&-1\\0&4&2\end{pmatrix}$
\item $\begin{pmatrix}2&-1&4\\0&3&-2\\1&0&5\end{pmatrix}$
\item $\begin{pmatrix}-1&0&2\\3&-5&1\\2&4&0\end{pmatrix}$
\item $\begin{pmatrix}1&3&-1\\-2&0&4\\0&5&2\end{pmatrix}$
\item $\begin{pmatrix}4&0&-3\\2&1&5\\-1&-2&0\end{pmatrix}$
\item $\begin{pmatrix}0&-3&1\\5&2&0\\-4&0&3\end{pmatrix}$
\item $\begin{pmatrix}3&1&0\\-2&4&-1\\0&-3&5\end{pmatrix}$
\item $\begin{pmatrix}-1&2&-3\\4&0&1\\0&-5&2\end{pmatrix}$
\item $\begin{pmatrix}2&0&-1\\0&3&4\\-3&1&0\end{pmatrix}$
\item $\begin{pmatrix}1&-4&0\\3&2&-5\\-2&0&1\end{pmatrix}$
\item $\begin{pmatrix}0&5&-2\\-1&0&3\\4&-3&0\end{pmatrix}$
\item $\begin{pmatrix}-2&0&4\\1&-3&0\\0&5&-1\end{pmatrix}$
\item $\begin{pmatrix}3&-2&0\\0&1&-4\\5&0&2\end{pmatrix}$
\item $\begin{pmatrix}1&0&-5\\-3&4&0\\2&-1&3\end{pmatrix}$
\item $\begin{pmatrix}0&-1&3\\2&5&0\\-4&0&1\end{pmatrix}$
\item $\begin{pmatrix}4&3&-1\\0&-2&5\\1&0&0\end{pmatrix}$
\item $\begin{pmatrix}-3&0&2\\1&4&0\\0&-1&5\end{pmatrix}$
\item $\begin{pmatrix}2&-5&1\\0&3&0\\-1&0&4\end{pmatrix}$
\item $\begin{pmatrix}0&2&-4\\-3&0&1\\5&-1&0\end{pmatrix}$
\item $\begin{pmatrix}1&1&-2\\-1&2&3\\3&0&-1\end{pmatrix}$
\item $\begin{pmatrix}-4&3&0\\2&0&-5\\0&1&3\end{pmatrix}$
\end{enumerate}
\end{multicols}
\emph{3. Solve the matrix equation $AX=B$.}
\bigskip
\begin{enumerate}[label=\arabic*.,  leftmargin=*]
\item $A=\begin{pmatrix}3&1\\5&2\end{pmatrix}$, $B=\begin{pmatrix}1\\-1\end{pmatrix}$.
\item $A=\begin{pmatrix}1&2\\3&5\end{pmatrix}$, $B=\begin{pmatrix}4&1\\2&-3\end{pmatrix}$.
\item $A=\begin{pmatrix}-2&3\\1&-1\end{pmatrix}$, $B=\begin{pmatrix}5\\-2\end{pmatrix}$.
\item $A=\begin{pmatrix}5&2\\-1&3\end{pmatrix}$, $B=\begin{pmatrix}3&0\\1&-1\end{pmatrix}$.
\item $A=\begin{pmatrix}1&-1\\-1&2\end{pmatrix}$, $B=\begin{pmatrix}2\\3\end{pmatrix}$.
\item $A=\begin{pmatrix}7&2\\3&1\end{pmatrix}$, $B=\begin{pmatrix}-1&4\\2&-3\end{pmatrix}$.
\item $A=\begin{pmatrix}-4&1\\2&3\end{pmatrix}$, $B=\begin{pmatrix}0\\5\end{pmatrix}$.
\item $A=\begin{pmatrix}2&1\\5&3\end{pmatrix}$, $B=\begin{pmatrix}1&-2\\3&4\end{pmatrix}$.
\item $A=\begin{pmatrix}8&5\\3&2\end{pmatrix}$, $B=\begin{pmatrix}2\\-1\end{pmatrix}$.
\item $A=\begin{pmatrix}9&-1\\4&1\end{pmatrix}$, $B=\begin{pmatrix}3&5\\0&-2\end{pmatrix}$.
\item $A=\begin{pmatrix}3&-5\\1&-2\end{pmatrix}$, $B=\begin{pmatrix}-1\\4\end{pmatrix}$.
\item $A=\begin{pmatrix}6&1\\-2&3\end{pmatrix}$, $B=\begin{pmatrix}2&-1\\0&5\end{pmatrix}$.
\item $A=\begin{pmatrix}1&4\\2&7\end{pmatrix}$, $B=\begin{pmatrix}3\\-2\end{pmatrix}$.
\item $A=\begin{pmatrix}0&1\\1&0\end{pmatrix}$, $B=\begin{pmatrix}5&-3\\1&2\end{pmatrix}$.
\item $A=\begin{pmatrix}-1&3\\2&-5\end{pmatrix}$, $B=\begin{pmatrix}4\\-7\end{pmatrix}$.
\item $A=\begin{pmatrix}4&3\\1&1\end{pmatrix}$, $B=\begin{pmatrix}2&0\\-1&3\end{pmatrix}$.
\item $A=\begin{pmatrix}2&-1\\-3&2\end{pmatrix}$, $B=\begin{pmatrix}0\\1\end{pmatrix}$.
\item $A=\begin{pmatrix}5&0\\0&-3\end{pmatrix}$, $B=\begin{pmatrix}10&-5\\6&-9\end{pmatrix}$.
\item $A=\begin{pmatrix}1&0\\0&1\end{pmatrix}$, $B=\begin{pmatrix}7\\-3\end{pmatrix}$.
\item $A=\begin{pmatrix}0&-2\\3&0\end{pmatrix}$, $B=\begin{pmatrix}4&-6\\9&3\end{pmatrix}$.
\item $A=\begin{pmatrix}7&-3\\2&1\end{pmatrix}$, $B=\begin{pmatrix}1\\0\end{pmatrix}$.
\item $A=\begin{pmatrix}-3&-1\\6&2\end{pmatrix}$, $B=\begin{pmatrix}2\\-4\end{pmatrix}$.
\item $A=\begin{pmatrix}2&3\\1&2\end{pmatrix}$, $B=\begin{pmatrix}1&-1\\0&2\end{pmatrix}$.
\item $A=\begin{pmatrix}1&-2\\3&-5\end{pmatrix}$, $B=\begin{pmatrix}3\\1\end{pmatrix}$.
\item $A=\begin{pmatrix}-1&-4\\3&11\end{pmatrix}$, $B=\begin{pmatrix}0&2\\-1&5\end{pmatrix}$.
\item $A=\begin{pmatrix}4&-1\\-2&1\end{pmatrix}$, $B=\begin{pmatrix}6\\-3\end{pmatrix}$.
\item $A=\begin{pmatrix}3&2\\4&3\end{pmatrix}$, $B=\begin{pmatrix}5&1\\-2&4\end{pmatrix}$.
\item $A=\begin{pmatrix}-5&2\\10&-4\end{pmatrix}$, $B=\begin{pmatrix}1\\3\end{pmatrix}$.
\item $A=\begin{pmatrix}1&1\\2&3\end{pmatrix}$, $B=\begin{pmatrix}-1&4\\2&-5\end{pmatrix}$.
\item $A=\begin{pmatrix}2&-1\\1&1\end{pmatrix}$, $B=\begin{pmatrix}3&0\\-2&1\end{pmatrix}$.
\end{enumerate}
\emph{4. Solve the matrix equation $XA=B$.}
\bigskip
\begin{enumerate}[label=\arabic*.,  leftmargin=*]
\item $A=\begin{pmatrix}1&2\\3&5\end{pmatrix}$, $B=\begin{pmatrix}4&-1\\2&3\end{pmatrix}$.
\item $A=\begin{pmatrix}3&1\\5&2\end{pmatrix}$, $B=\begin{pmatrix}-2&1\\0&3\end{pmatrix}$.
\item $A=\begin{pmatrix}-1&3\\2&-5\end{pmatrix}$, $B=\begin{pmatrix}1&0\\-3&2\end{pmatrix}$.
\item $A=\begin{pmatrix}2&1\\5&3\end{pmatrix}$, $B=\begin{pmatrix}0&-1\\4&2\end{pmatrix}$.
\item $A=\begin{pmatrix}7&2\\3&1\end{pmatrix}$, $B=\begin{pmatrix}5&1\\-2&0\end{pmatrix}$.
\item $A=\begin{pmatrix}1&-1\\-1&2\end{pmatrix}$, $B=\begin{pmatrix}3&-4\\1&0\end{pmatrix}$.
\item $A=\begin{pmatrix}5&2\\-1&3\end{pmatrix}$, $B=\begin{pmatrix}-1&2\\3&-5\end{pmatrix}$.
\item $A=\begin{pmatrix}-2&3\\1&-1\end{pmatrix}$, $B=\begin{pmatrix}0&1\\-2&4\end{pmatrix}$.
\item $A=\begin{pmatrix}8&5\\3&2\end{pmatrix}$, $B=\begin{pmatrix}1&-1\\4&3\end{pmatrix}$.
\item $A=\begin{pmatrix}0&1\\1&0\end{pmatrix}$, $B=\begin{pmatrix}2&-3\\5&1\end{pmatrix}$.
\item $A=\begin{pmatrix}4&-1\\2&3\end{pmatrix}$, $B=\begin{pmatrix}6&0\\-1&5\end{pmatrix}$.
\item $A=\begin{pmatrix}9&-1\\4&1\end{pmatrix}$, $B=\begin{pmatrix}2&-2\\1&4\end{pmatrix}$.
\item $A=\begin{pmatrix}3&-5\\1&-2\end{pmatrix}$, $B=\begin{pmatrix}0&3\\-1&2\end{pmatrix}$.
\item $A=\begin{pmatrix}6&1\\-2&3\end{pmatrix}$, $B=\begin{pmatrix}-3&1\\2&0\end{pmatrix}$.
\item $A=\begin{pmatrix}1&4\\2&7\end{pmatrix}$, $B=\begin{pmatrix}3&-1\\0&2\end{pmatrix}$.
\item $A=\begin{pmatrix}-4&1\\2&3\end{pmatrix}$, $B=\begin{pmatrix}1&-1\\5&2\end{pmatrix}$.
\item $A=\begin{pmatrix}2&-1\\-3&2\end{pmatrix}$, $B=\begin{pmatrix}4&0\\-2&3\end{pmatrix}$.
\item $A=\begin{pmatrix}5&0\\0&-3\end{pmatrix}$, $B=\begin{pmatrix}10&0\\0&9\end{pmatrix}$.
\item $A=\begin{pmatrix}3&2\\4&3\end{pmatrix}$, $B=\begin{pmatrix}1&-2\\3&0\end{pmatrix}$.
\item $A=\begin{pmatrix}0&-2\\3&0\end{pmatrix}$, $B=\begin{pmatrix}6&4\\-3&6\end{pmatrix}$.
\item $A=\begin{pmatrix}7&-3\\2&1\end{pmatrix}$, $B=\begin{pmatrix}0&-1\\2&5\end{pmatrix}$.
\item $A=\begin{pmatrix}1&-2\\3&-5\end{pmatrix}$, $B=\begin{pmatrix}2&0\\-3&1\end{pmatrix}$.
\item $A=\begin{pmatrix}2&3\\1&2\end{pmatrix}$, $B=\begin{pmatrix}-1&3\\4&-2\end{pmatrix}$.
\item $A=\begin{pmatrix}-1&-4\\3&11\end{pmatrix}$, $B=\begin{pmatrix}5&0\\-1&3\end{pmatrix}$.
\item $A=\begin{pmatrix}4&3\\1&1\end{pmatrix}$, $B=\begin{pmatrix}0&1\\-2&4\end{pmatrix}$.
\item $A=\begin{pmatrix}1&1\\2&3\end{pmatrix}$, $B=\begin{pmatrix}4&-3\\2&1\end{pmatrix}$.
\item $A=\begin{pmatrix}2&-1\\1&1\end{pmatrix}$, $B=\begin{pmatrix}0&3\\-1&2\end{pmatrix}$.
\item $A=\begin{pmatrix}-3&2\\1&0\end{pmatrix}$, $B=\begin{pmatrix}6&-4\\3&1\end{pmatrix}$.
\item $A=\begin{pmatrix}5&-2\\-3&1\end{pmatrix}$, $B=\begin{pmatrix}1&1\\-2&0\end{pmatrix}$.
\item $A=\begin{pmatrix}4&-1\\-2&1\end{pmatrix}$, $B=\begin{pmatrix}2&-3\\0&5\end{pmatrix}$.
\end{enumerate}
